\section{对称形式幂级数}

记号约定:$\mathbf{X}\coloneq (X_{i})_{i\in I},\mathbf{S}\coloneq\left(S_{i}\right)_{i=1}^{\infty}$是不定元族.

\begin{proposition}{}{}
	对每个$\sigma\in\mathfrak{S}_{I}$,存在唯一的连续自同构$\varphi_{\sigma}:A[[\mathbf{X}]]\to A[[\mathbf{X}]]$,对每个$i\in I$有$\varphi_{\sigma}\left(X_{i}\right)=X_{\sigma\left(i\right)}$.
\end{proposition}

\begin{definition}{对称形式幂级数}{}
	设$f\in A[[\mathbf{X}]]$.若对每个$\sigma\in\mathfrak{S}_{I}$有$\varphi_{\sigma}\left(f\right)=f$,则称$f$是对称形式幂级数.记
	\begin{align*}
		A[[\mathbf{X}]]^{\mathfrak{S}_{I}}\coloneq\left\{g\in A[[\mathbf{X}]]\middle|\forall\sigma\in\mathfrak{S}_{I}\left(\varphi_{\sigma}\left(g\right)=g\right)\right\},
	\end{align*}
	并配备$A[[\mathbf{X}]]$诱导的子空间拓扑,那么这是$A[[\mathbf{X}]]$的闭子代数.
\end{definition}

\begin{lemma}{广义初等对称函数的良定义性}{}
	在$A[[\mathbf{X},T]]$中,$\left(1+X_{i}T\right)_{i\in I}$是可乘的.
\end{lemma}

\begin{definition}{广义初等对称形式幂级数}{}
	对每个$k>0$,称$s_{k}\coloneq\sum\limits_{H\in\fin_{k}\left(I\right)}\left(\prod\limits_{i\in H}X_{i}\right)\in A[[\mathbf{X}]]$是$k$次初等对称形式幂级数.
\end{definition}

\begin{proposition}{}{}
	存在唯一的连续$A$-代数同态$\varphi_{I}:A[[\mathbf{S}]]\to A[[\mathbf{X}]]^{\mathfrak{S}_{I}}$使得$\forall k\geq 1\left(\varphi_{I}\left(S_{k}\right)=s_{k}\right)$.
\end{proposition}

\begin{lemma}{形式幂级数的截断消没引理}{}
	设$J\in\fin\left(I\right),0<r\leq\left|J\right|$,$f\in A[[\mathbf{X}]]^{\mathfrak{S}_{I}}$且全次数为$r$.记$f_{0}$是把$f$中的不定元$X_{i}\left(i\in I\setminus J\right)$换成$0$所得的形式幂级数,则
	\begin{align*}
		f_{0}=0\implies f=0.
	\end{align*}
\end{lemma}

\begin{lemma}{齐次对称形式幂级数的提升}{}
	设$f\in A[[\mathbf{S}]]^{\mathfrak{S}_{I}}$且全次数为$r$.为$A[\mathbf{S}]$配备$\N$型分次,使得对每个$k\geq 1$,多项式$X_{k}$的权是$k$.那么,存在唯一的总权为$r$的等权多项式$P\in A[\mathbf{S}]$使得$f=\psi_{I}\left(P\right)$.
\end{lemma}

\begin{lemma}{对称形式幂级数的邻域基}{}
	对每个$m\in\N$,记$\mathfrak{c}_{m}=\left\{f\in A[[\mathbf{X}]]^{\mathfrak{S}_{I}}\middle|\omega\left(f\right)\geq m\right\}$,则$\left(\mathfrak{c}_{m}\right)_{m\in\N}$是$A[[\mathbf{X}]]^{\mathfrak{S}_{I}}$中$0$处的邻域基.
\end{lemma}


\begin{theorem}{对称形式幂级数基本定理}{}
	\begin{enumerate}
		\item 若$\left|I\right|=n\in\N$,则$\varphi_{I}$诱导了拓扑$A$-代数同构$A[[S_{1},\ldots,S_{n}]]\to A[[\mathbf{X}]]^{\mathfrak{S}_{I}}$.
		\item 若$I$无限,则$\varphi_{I}$诱导了拓扑$A$-代数同构$A[[\mathbf{S}]]\to A[[\mathbf{X}]]^{\mathfrak{S}_{I}}$.
	\end{enumerate}
\end{theorem}

















































































